% ================================================================
%  练习 (原书书末 Exercises, 书页 295--298)
%  原书为不编号的章级标题, 紧随第 8 章之后; 公式 (8.97) 仍按第 8 章编号,
%  故本节用 \chapter* 排, 不推进 chapter 计数器.
% ================================================================
\chapter*{练习}
\addcontentsline{toc}{chapter}{练习}
\markboth{练习}{练习}

{\setlength{\emergencystretch}{2em}%
本节的练习按原书顺序编号为 Ex 1--Ex 35, 覆盖全书各章. 少数题目需用到
前面各章的结论, 均以交叉引用给出.\par}

\begin{enumerate}[label=\textbf{Ex \arabic*.},leftmargin=3.7em,
  labelsep=0.5em,itemsep=0.85em,topsep=0.5em,parsep=0.35em]

\item 证明图上的面积泛函是凸的.

\item 证明悬链面与螺旋面都是极小曲面.

\item 验证 Weierstrass 表示
\[
F(z)=\operatorname{Re}\int_{\zeta\in\gamma_{z_0,z}}
\Bigl(\tfrac12\bigl(g^{-1}(\zeta)-g(\zeta)\bigr),\,
\tfrac{\ii}{2}\bigl(g^{-1}(\zeta)+g(\zeta)\bigr),\,1\Bigr)\phi(\zeta),
\]
其中 $g(z)=z$, $\phi(z)=dz/z$, $\Omega=\C\setminus\{0\}$, 给出一个悬链面.

\item 由极小曲面的 Weierstrass 数据推导其单位法向量、度量与曲率的公式
\eqref{eq:1.115}--\eqref{eq:1.117}.

\item 计算悬链面的总曲率, 即 $\int K$.

\item 求极小曲面方程在 $\R^2\setminus D_1$ 上的所有旋转对称解.

\item 证明 Enneper 曲面不是嵌入的, 但其单位法向量渐近竖直(即像一张图).
Enneper 曲面渐近地是什么样子?

\item 设 $M$ 是截面曲率至多为 $-1$ 的三维流形.
  \begin{itemize}[label=\textbullet,leftmargin=1.6em,itemsep=0.2em,topsep=0.3em]
  \item 证明不存在 $S^2$ 到 $M$ 的极小浸入.
  \item 对亏格 $g$ 的曲面到 $M$ 的极小浸入, 证明一个面积界.
  \end{itemize}

\item 设 $u\colon D_1\to\R$ 满足极小曲面方程, 且在 $D_1$ 上 $|Du|\leq1$.
用标准椭圆理论证明一致界 $|D^2u|(0)\leq C$.

\item {\setlength{\emergencystretch}{2em}%
求一列函数 $u_j\colon D_1\to\R$, 它们满足极小曲面方程, 在 $D_1$ 上
$|Du_j|\leq1$, 但 $\sup_{D_1}|D^2u_j|\to\infty$.\par}

\item 证明推论~\ref{cor:1.8} 中定义的通量同态决定了悬链面的定向(即旋转轴的
方向)与尺度.

\item 补全引理~\ref{lem:1.18} 证明的细节.

\item 证明悬链面上的任何正调和函数必为常数.

\item 设 $\Sigma$ 是极小环带, 其边界分支为圆
$\{x^2+y^2=1,\ z=H\}$ 与 $\{x^2+y^2=1,\ z=-H\}$. 证明 $H$ 不能太大.

\item 设 $\Sigma$ 是完全极小曲面, 且对所有的 $R$ 有
\[
\Area\bigl(B_R\cap\Sigma\bigr)\leq2\pi R^2 .
\]
证明 $\Sigma$ 要么是嵌入的, 要么是两张平面的并.

\item 对非负函数 $f$ 定义散度型算子
\[
Lu=\dive\bigl(f(|\nabla u|^2)\nabla u\bigr)
\]
(例如 $f(s)=(1+s)^{-1/2}$ 给出极小曲面方程). 问: 对哪些函数 $f$, $L$
具有强极值原理?

\item 通过证明极小图的唯一性来补全定理~\ref{thm:1.29} 的证明.

\item 悬链面的 Morse 指标是多少?

\item 证明螺旋面的 Morse 指标是无穷. 你能把螺旋面分成两块稳定部分吗?

\item 证明引理~\ref{lem:3.24}.

\item Enneper 曲面的 Morse 指标是多少?

\item 分类 $\R^3$ 中所有具有一一高斯映射的完全浸入极小曲面.

\item 证明引理~\ref{lem:1.18}.

\item 设 $\Sigma$ 是 $\R^3$ 中的单位球面. 用 \eqref{eq:2.4} 计算曲率张量
$R_{ijk\ell}$.

\item 设 $\Sigma\subset\R^3$ 是极小曲面. 共形度量 $|A|^2\langle\cdot,\cdot\rangle$
的曲率是多少?

\item 设 $\Sigma$ 是极小图, 令 $v=(1+|\nabla_{\R^2}u|^2)^{1/2}$(即单位法向量
第三个分量的倒数). 计算
\[
\Delta_\Sigma\bigl(|A|^2v^2\bigr).
\]
并用它由极值原理估计 $|A|^2(0)$.

\item 用尺度化证明下述``小能量估计'':

存在 $\epsilon>0$, 使得若在 $B_R\subset\R^3$ 上 $\Delta u=-u^5$ 且
\[
\int_{B_R}u^6<\epsilon,
\]
则 $|u(0)|\leq R^{-1/2}$.

\item 改用积分方法(即分部积分与 Sobolev 不等式)重做上一题, 而不使用尺度化.

\item 证明存在 $\epsilon>0$, 使得若 $u\colon\R^2\to\bS^2\subset\R^3$ 满足
``调和映射方程''
\[
\Delta u=-|\nabla u|^2u
\]
且能量很小
\[
\int|\nabla u|^2<\epsilon,
\]
则 $u$ 为常数.

\item 设 $\Sigma^2\subset\bS^3$ 是浸入的极小(拓扑)$\bS^2$. 证明 $\Sigma$ 是
全测地的(从而是``圆''的). 提示: 证明 $A$ 可以看作一个全纯二次微分.

\item 证明 Heinz 估计的下述加强形式:

存在常数 $C$, 使得若 $u$ 的图在 $D_1\subset\R^2$ 上是极小的, 则在(其上方)
点 $0$ 处的曲率 $|A|^2$ 满足
\[
|A|^2\leq \frac{C}{1+|\nabla u(0)|^2}.
\]

\item 证明 $\bS^3$ 中闭极小曲面的最小可能面积是 $4\pi$.

\end{enumerate}

{\setlength{\emergencystretch}{2em}%
对下面三个练习, 在 $\R^{n+1}$ 的超曲面 $\Sigma$ 上定义泛函 $F$:
\begin{equation}
F(\Sigma)=\int_\Sigma \exp\Bigl(-\frac{|x|^2}{4}\Bigr).
\label{eq:8.97}
\end{equation}
\par}

\begin{enumerate}[label=\textbf{Ex \arabic*.},leftmargin=3.7em,
  labelsep=0.5em,itemsep=0.85em,topsep=0.5em,parsep=0.35em,start=33]

\item 计算 $F$ 的第一变分以及 $F$ 的临界点的 Euler--Lagrange 方程.
(这些临界点称为\newterm{self-shrinker}{自收缩子}{self-shrinker}, 它们出现在
平均曲率流中; 详见 \cite{Hu}、\cite{CM25} 与 \cite{CM6}.)

\item 对紧支撑的法向变分, 计算 $F$ 的第二变分.

\item 证明不存在作为 $F$ 的稳定临界点的闭超曲面.

\end{enumerate}
